% Assignment source: ne630/homework/html/HW09.html.
% Legacy foundation: ne630_problems/lesson_09/statement.tex.

\noindent{\bf Problem 1}. Verify Eqs.~(3.23) and (3.25):
\[
 \Sigma_s(E)\varphi(E)=\frac{C}{E}
 \tag{FNRP 3.23}
\]
and
\[
 q=\left[1+\frac{\alpha}{1-\alpha}\ln\alpha\right]C.
 \tag{FNRP 3.25}
\]
\textbf{Hint:} To verify Eq.~(3.23), substitute Eq.~(3.23) into Eq.~(3.22)
and show that the left- and right-hand sides are equal.

\begin{adaptedsolution}
Substitute $\Sigma_s(E)\phi(E)=C/E$ into Eq.~(3.22):
\begin{align*}
 \Sigma_s(E)\phi(E)
 &=\int_E^{E/\alpha}
   \frac{1}{(1-\alpha)E'}\Sigma_s(E')\phi(E')\,dE'\\
 &=\int_E^{E/\alpha}\frac{1}{(1-\alpha)E'}\frac{C}{E'}\,dE'\\
 &=\frac{C}{1-\alpha}\int_E^{E/\alpha}\frac{dE'}{E'^2}\\
 &=\frac{C}{1-\alpha}\left[-\frac{1}{E'}\right]_E^{E/\alpha}\\
 &=\frac{C}{1-\alpha}\left(\frac{1}{E}-\frac{\alpha}{E}\right)
  =\boxed{\frac{C}{E}}.
\end{align*}
For Eq.~(3.25),
\begin{align*}
 q
 &=\int_E^{E/\alpha}
   \left[\int_{\alpha E'}^E
   \frac{1}{(1-\alpha)E'}\Sigma_s(E')\phi(E')\,dE''\right]dE'\\
 &=\frac{C}{1-\alpha}\int_E^{E/\alpha}
   \frac{E-\alpha E'}{E'^2}\,dE'\\
 &=\frac{C}{1-\alpha}
   \left(\int_E^{E/\alpha}\frac{E}{E'^2}\,dE'
         -\int_E^{E/\alpha}\frac{\alpha}{E'}\,dE'\right)\\
 &=\frac{C}{1-\alpha}\left(1-\alpha+\alpha\ln\alpha\right)\\
 &=\boxed{\left[1+\frac{\alpha}{1-\alpha}\ln\alpha\right]C}.
\end{align*}
\end{adaptedsolution}

\newpage

\noindent{\bf Problem 2}. For thermal neutrons, calculate $\bar{\eta}$ as a
function of uranium enrichment and plot your results for enrichments between
0\% and 30\%. Use the uranium data from the following table:
\[
\begin{array}{lrrr}
\toprule
 & \nu & \sigma_f~(\mathrm{barns}) & \sigma_a~(\mathrm{barns})\\
\midrule
\text{Uranium-235} & 2.43 & 505 & 591\\
\text{Plutonium-239} & 2.90 & 698 & 973\\
\text{Uranium-238} & \text{---} & 0 & 2.42\\
\bottomrule
\end{array}
\]
Consider the value you compute for 4\% enrichment. How does this value compare
with the values computed for Problem 2 of HW08?

\begin{adaptedsolution}
For enrichment $e$, the thermal reproduction factor for the uranium mixture is
\[
 \eta(e)
 =\frac{e\,\nu_{235}\sigma_{f,235}}
        {e\,\sigma_{a,235}+(1-e)\sigma_{a,238}}.
\]
Substitution gives
\[
 \eta(e)=
 \frac{e(2.43)(505)}
      {e(591)+(1-e)(2.42)}.
\]
Thus, for $e=0.04$,
\[
 \eta(0.04)
 =\frac{0.04(2.43)(505)}
      {0.04(591)+0.96(2.42)}
 \approx \boxed{1.89}.
\]
A Python plotting solution is
\begin{minted}{python}
import numpy as np
import matplotlib.pyplot as plt

nu_235 = 2.43
sigma_f_235 = 505.0
sigma_a_235 = 591.0
sigma_a_238 = 2.42

enrichment = np.linspace(0, 0.3, 100)
eta = lambda e: e*nu_235*sigma_f_235/(e*sigma_a_235
                                      + (1-e)*sigma_a_238)
plt.plot(100*enrichment, eta(enrichment), "k")
plt.plot(4, eta(0.04), "rx")
plt.xlabel("enrichment (%)")
plt.ylabel(r"$\eta$")
\end{minted}
The 4\% value is close to the thermal-range values from HW08 and to the low
energy portion of the 4\% $\eta(E)$ curve.  It is much more favorable than the
resonance-dominated epithermal behavior.
\end{adaptedsolution}

\newpage

\noindent{\bf Problem 3}. Lethargy, defined as $u=\ln(E_0/E)$, is often used
in neutron slowing-down problems; lethargy increases as energy decreases. Note
the following transformations:
\[
 \varphi(E)\,dE=-\varphi(u)\,du,\qquad
 p(E\to E')\,dE'=-p(u\to u')\,du',\qquad
 \Sigma_x(E)=\Sigma_x(u).
\]
\begin{enumerate}[label=(\alph*)]
\item Show that $p(E\to E')$, given by Eq.~(2.47), becomes
\[
p(u\to u')=
\begin{cases}
\displaystyle\frac{1}{1-\alpha}\exp(u-u'),
  & u\le u'\le u+\ln(1/\alpha),\\
0, & \text{otherwise}.
\end{cases}
\]
\item Express Eq.~(3.22) in terms of $u$.
\item Show that the slowing-down decrement can be expressed as
$\xi=\overline{\Delta u}$.
\end{enumerate}

\begin{adaptedsolution}
\begin{enumerate}[label=(\alph*)]
\item Since $E'=E_0e^{-u'}$, $dE'/du'=-E'$.  Therefore,
\[
 p(u\to u')
 =-p(E\to E')\frac{dE'}{du'}
 =\frac{1}{(1-\alpha)E}E'
 =\frac{1}{1-\alpha}\exp(u-u').
\]
The original support $\alpha E\le E'\le E$ becomes
$u\le u'\le u+\ln(1/\alpha)$.
\item Using $\phi(u)=\phi(E)E$ and $E=E_0e^{-u}$, Eq.~(3.22) becomes
\[
 \boxed{\Sigma_s(u)\phi(u)
 =\int_{u-\ln(1/\alpha)}^{u}
  \frac{e^{u'-u}}{1-\alpha}\Sigma_s(u')\phi(u')\,du'}.
\]
\item With the density found in part (a),
\begin{align*}
 \overline{\Delta u}
 &=\int_u^{u+\ln(1/\alpha)}
   \frac{1}{1-\alpha}e^{u-u'}(u'-u)\,du'\\
 &=1+\frac{\alpha}{1-\alpha}\ln\alpha\\
 &=\boxed{\xi}.
\end{align*}
\end{enumerate}
\end{adaptedsolution}

\clearpage
